\section{Restricting to the edges incident with a vertex block}
\label{sec:graph}

Let $E_n=\binom{[n]}2$ be the edge set of the complete graph on $[n]$,
and let $f\colon\{0,1\}^{E_n}\to\{0,1\}$ be invariant under vertex
permutations.  Fix $p\in(0,1)$ and put $\sigma=\sqrt{p(1-p)}$.
All Fourier coefficients, variances, and influences in this section
use this bias.  Monotonicity is not needed.

The purpose of a vertex block is to connect two different notions of
degree.  A restriction to the edges incident with the block retains
enough symmetry to apply Lemma~\ref{lem:low-degree}.  If the block meets
the nonisolated vertices of a Fourier support at just one vertex $u$,
then the surviving Fourier degree is the graph degree of $u$.
An $s$-edge graph with $s\ge1$ has a nonisolated vertex of degree at most
$\sqrt{2s}$.  Thus a bound for conditional Fourier degrees up to $k$
will control original degrees up to $k^2/2$.

The restriction identity is the biased-product form of the standard
Fourier restriction identity; compare
\cite[arXiv version~2, Lemma~2.7]{KKLMS20}.
We include its short proof. The next identities hold for every Boolean
$f$, without a symmetry assumption.

\begin{lemma}[Exact averaging over a restriction]
\label{lem:restriction}
Let $F\subseteq E_n$, put $H=E_n\setminus F$, and, for
$y\in\{0,1\}^{H}$, let $f_y(x)=f(x,y)$ for
$x\in\{0,1\}^{F}$.  If $y$ has the Bernoulli-$p$ product law, then,
for every real $k>0$,
\begin{align}
 \E_y\sum_{\substack{A\subseteq F\\1\le |A|\le k}}
        \widehat{f_y}(A)^2
 &=\sum_{\substack{S\subseteq E_n\\1\le|S\cap F|\le k}}
        \widehat f(S)^2,
 \label{eq:restriction-mass}\\
 \E_y I(f_y)&=\sum_{e\in F}I_e(f).
 \label{eq:restriction-influence}
\end{align}
\end{lemma}

\begin{proof}
Expand $f$ in the product Fourier basis and integrate over $x$.
For each $A\subseteq F$ this gives the exact finite identity
\[
 \widehat{f_y}(A)
 =\sum_{D\subseteq H}\widehat f(A\cup D)\chi_D(y).
\]
Orthonormality on the exterior coordinates therefore gives
$\E_y\widehat{f_y}(A)^2
=\sum_{D\subseteq H}\widehat f(A\cup D)^2$.
Summing over $A$ proves \eqref{eq:restriction-mass}.  In particular,
there is no constraint on the exterior Fourier degree $|D|$.

For $e\in F$, the pivotal event for $f_y$ at $e$ is the pivotal event
for $f$ at $e$ with the exterior coordinates fixed to $y$.  Averaging
over $y$ and the remaining free coordinates gives
$\E_y I_e(f_y)=I_e(f)$; summation proves
\eqref{eq:restriction-influence}.
\end{proof}

\begin{proposition}[The block estimate]
\label{prop:block}
Let $3\le m\le n$ be an integer, let $B$ be uniform among the
$m$-element subsets of $[n]$, and set
\[
 F_B=\{e\in E_n:e\cap B\ne\varnothing\},
 \qquad k=\frac{\sigma\log m}{64}.
\]
Then
\begin{equation}
 \sum_{S\subseteq E_n}
 \Pr_B\{1\le |S\cap F_B|\le k\}\widehat f(S)^2
 \le \frac{2m}{n}\,m^{-1/4}I(f).
 \label{eq:block-average}
\end{equation}
\end{proposition}

\begin{proof}
Fix $B$, and fix an arbitrary assignment $y$ to
$H_B=E_n\setminus F_B$.  The free set $F_B$ contains \emph{all} edges
with at least one endpoint in $B$: both edges within $B$ and edges
from $B$ to its complement.  Every permutation supported on $B$
fixes each edge of $H_B$ individually, and hence fixes $y$.
The restricted function $f_y$ is consequently invariant under this
permutation group even if the graph encoded by $y$ has no symmetry.

The free-edge orbits are precisely the $\binom m2$ edges within $B$
and, for each \emph{fixed} $w\in[n]\setminus B$, the $m$ edges
$\{\{u,w\}:u\in B\}$.  Every orbit has size at least $m$, since
$m\ge3$.  The product law and the pivotal events are preserved by
these permutations, so influences agree within each orbit.
It follows that
\[
 I_e(f_y)\le\frac{I(f_y)}m\qquad(e\in F_B).
\]
Lemma~\ref{lem:low-degree} applies separately to every $y$.  Its
right-hand side is linear in $I(f_y)$, so averaging and using
Lemma~\ref{lem:restriction} gives
\[
 \sum_{\substack{S\subseteq E_n\\1\le |S\cap F_B|\le k}}
       \widehat f(S)^2
 \le m^{-1/4}\sum_{e\in F_B}I_e(f).
\]
Finally average over $B$.  Each endpoint of a fixed edge belongs to
$B$ with probability $m/n$, whence
$\Pr_B(e\in F_B)\le2m/n$.  This proves
\eqref{eq:block-average}.
\end{proof}

\begin{lemma}[Capturing low Fourier levels]
\label{lem:block-capture}
Suppose $t=\log n\ge16$, and put
\[
 m=\lfloor\sqrt n\rfloor,
 \qquad k=\frac{\sigma\log m}{64}.
\]
Then
\begin{equation}
 \sum_{\substack{S\subseteq E_n\\1\le |S|\le k^2/2}}
       \widehat f(S)^2
 \le 4m^{-1/4}I(f).
 \label{eq:captured-mass}
\end{equation}
\end{lemma}

\begin{proof}
Fix a nonempty $S\subseteq E_n$ with $s=|S|\le k^2/2$, viewed as
the edge set of a simple graph.  Let $V(S)$ be its set of nonisolated
vertices and write $v=|V(S)|$.  Since $s\le\binom v2<v^2/2$, its
average degree over these vertices satisfies
\[
 \frac{2s}{v}<\sqrt{2s}.
\]
Choose a vertex $u\in V(S)$ whose degree $d_S(u)$ is at most this average.  Then
\begin{equation}
 1\le d_S(u)\le\sqrt{2s}\le k,
 \qquad v\le2s\le k^2.
 \label{eq:support-degree}
\end{equation}

Choose $B$ uniformly among the $m$-element vertex sets.  Conditioned
on $u\in B$, each other specified vertex belongs to $B$ with
probability $(m-1)/(n-1)$.  A union bound gives
\begin{equation}
 \Pr_B\{B\cap V(S)=\{u\}\}
 \ge \frac mn
 \left(1-(v-1)\frac{m-1}{n-1}\right).
 \label{eq:one-vertex-event}
\end{equation}
On this event, $S\cap F_B$ consists of \emph{every} support edge
incident with $u$, and contains no other support edge.  Its size is
therefore $d_S(u)$, between $1$ and $k$; it need not be $1$.
Figure~\ref{fig:block} illustrates this distinction.

\begin{figure}[tb]
\centering
\setlength{\unitlength}{0.93pt}
\begin{picture}(460,190)
 \put(12,176){\makebox(202,12){\small (a) The free coordinate set $F_B$}}
 \put(246,176){\makebox(202,12){\small (b) One Fourier support $S$}}
 \linethickness{0.45pt}
 \put(40,164){\line(1,0){52}}
 \qbezier(92,164)(108,164)(108,148)
 \put(108,56){\line(0,1){92}}
 \qbezier(108,56)(108,40)(92,40)
 \put(40,40){\line(1,0){52}}
 \qbezier(40,40)(24,40)(24,56)
 \put(24,56){\line(0,1){92}}
 \qbezier(24,148)(24,164)(40,164)
 \put(274,164){\line(1,0){52}}
 \qbezier(326,164)(342,164)(342,148)
 \put(342,56){\line(0,1){92}}
 \qbezier(342,56)(342,40)(326,40)
 \put(274,40){\line(1,0){52}}
 \qbezier(274,40)(258,40)(258,56)
 \put(258,56){\line(0,1){92}}
 \qbezier(258,148)(258,164)(274,164)
 \put(23,151){\makebox(15,10){$B$}}
 \put(257,151){\makebox(15,10){$B$}}
 \thinlines
 \qbezier(47,133)(65.5,117)(84,101)
 \qbezier(84,101)(65.5,85)(47,69)
 \put(47,69){\line(0,1){64}}
 \qbezier(47,133)(110.5,142)(174,151)
 \qbezier(47,133)(121,125)(195,117)
 \qbezier(47,133)(121,103.5)(195,74)
 \qbezier(47,133)(110.5,86.5)(174,40)
 \qbezier(84,101)(129,126)(174,151)
 \qbezier(84,101)(139.5,109)(195,117)
 \qbezier(84,101)(139.5,87.5)(195,74)
 \qbezier(84,101)(129,70.5)(174,40)
 \qbezier(47,69)(110.5,110)(174,151)
 \qbezier(47,69)(121,93)(195,117)
 \qbezier(47,69)(121,71.5)(195,74)
 \qbezier(47,69)(110.5,54.5)(174,40)
 \put(47,133){\circle*{5}}
 \put(84,101){\circle*{5}}
 \put(47,69){\circle*{5}}
 \put(174,151){\circle*{5}}
 \put(195,117){\circle*{5}}
 \put(195,74){\circle*{5}}
 \put(174,40){\circle*{5}}
 \put(24,9){\makebox(190,22){\shortstack{\scriptsize All displayed edges are free.\\[-1pt]
 \scriptsize Edges wholly outside $B$ are fixed.}}}
 \linethickness{1pt}
 \qbezier(318,101)(363,126)(408,151)
 \qbezier(318,101)(373.5,109)(429,117)
 \qbezier(408,151)(418.5,134)(429,117)
 \put(281,133){\circle{5}}
 \put(318,101){\circle*{6}}
 \put(281,69){\circle{5}}
 \put(408,151){\circle*{6}}
 \put(429,117){\circle*{6}}
 \put(429,74){\circle{5}}
 \put(408,40){\circle{5}}
 \put(302,88){\makebox(13,12){$u$}}
 \put(410,156){\makebox(19,10){\scriptsize $v_1$}}
 \put(435,115){\makebox(19,10){\scriptsize $v_2$}}
 \put(435,72){\makebox(19,10){\scriptsize $v_3$}}
 \put(410,26){\makebox(19,10){\scriptsize $v_4$}}
 \put(327,143){\makebox(63,10){\scriptsize survives}}
 \put(430,146){\makebox(28,20){\shortstack{\scriptsize fixed\\[-1pt]\scriptsize by $y$}}}
 \thinlines
 \put(440,142){\vector(-1,-1){16}}
 \put(257,9){\makebox(190,15){\small $|S\cap F_B|=d_S(u)=2$}}
\end{picture}
\caption{The block restriction and the one-vertex event.
In (a), $F_B$ contains every edge incident with $B$, including its
internal edges; only edges with both endpoints outside $B$ are fixed.
In (b), only support edges are drawn, and filled vertices are the
nonisolated vertices $V(S)$.  The event $B\cap V(S)=\{u\}$ leaves
both support edges incident with $u$ free.  The exterior support
edge is fixed, and its contribution is retained by the exact
averaging identity \eqref{eq:restriction-mass}.}
\label{fig:block}
\end{figure}


To bound the probability uniformly in $p$, note that $k\le t$ and
\[
 \frac{m-1}{n-1}\le\frac1{\sqrt n},
 \qquad
 (v-1)\frac{m-1}{n-1}
 \le t^2e^{-t/2}\le256e^{-8}<\frac12.
\]
Here $t^2e^{-t/2}$ is decreasing for $t\ge16$.
Consequently, for every support under consideration,
\begin{equation}
 \Pr_B\{1\le|S\cap F_B|\le k\}\ge\frac{m}{2n}.
 \label{eq:support-detection}
\end{equation}
The choice of $u$ may depend on $S$.  The lower bound
\eqref{eq:support-detection} is applied to each nonnegative summand
$\widehat f(S)^2$ separately; no block must capture all supports
simultaneously.  Comparing with \eqref{eq:block-average} and
cancelling $m/(2n)$ proves \eqref{eq:captured-mass}.
If $k^2/2<1$, the sum in \eqref{eq:captured-mass} is empty and the
conclusion holds without choosing a support.
\end{proof}

\begin{proof}[Proof of Theorem~\ref{thm:variance}]
First suppose $t=\log n\ge16$, and choose $m$ and $k$ as in
Lemma~\ref{lem:block-capture}; in particular $m\ge3$ and $k>0$.
For the remaining Fourier degrees, the Dirichlet identity
\eqref{eq:dirichlet} gives
\begin{equation}
 \sum_{|S|>k^2/2}\widehat f(S)^2
 \le\frac2{k^2}\sum_S|S|\widehat f(S)^2
 =\frac{2\sigma^2}{k^2}I(f).
 \label{eq:graph-tail}
\end{equation}
This remains valid for $k^2/2<1$, when the left side includes all
nonconstant coefficients.  Combining \eqref{eq:parseval},
\eqref{eq:captured-mass}, and \eqref{eq:graph-tail} yields
\begin{equation}
 \Var(f)\le
 \left(4m^{-1/4}+\frac{2\sigma^2}{k^2}\right)I(f).
 \label{eq:graph-variance-preconstant}
\end{equation}

For completeness, all constants can be bounded independently of
$p$ and $n$.  Since $m\ge\sqrt n/2$ and $t\ge16$,
$\log m\ge t/2-\log2\ge t/3$.  Thus
\[
 \frac{2\sigma^2}{k^2}
 =\frac{8192}{(\log m)^2}\le\frac{73728}{t^2}.
\]
The cancellation of $\sigma^2$ is exact, even when $p$ is arbitrarily
close to $0$ or $1$.  Also, $t^2e^{-t/8}$ is nonincreasing for
$t\ge16$, so
\[
 4m^{-1/4}t^2
 \le4\,2^{1/4}t^2e^{-t/8}
 \le4\,2^{1/4}\,256e^{-2}<256.
\]
Thus the right-hand side of \eqref{eq:graph-variance-preconstant} is at most
$(73728+256)I(f)/t^2$, which is at most
$2^{17}I(f)/t^2$.

If instead $2\le n<e^{16}$, the same Fourier identities immediately
give
\[
 \Var(f)\le\sigma^2I(f)\le\frac14 I(f)
 \le\frac{64}{(\log n)^2}I(f).
\]
Hence the single choice $C_0=2^{17}$ works for every integer $n\ge2$
and every $p\in(0,1)$, as asserted.
\end{proof}
